\documentclass[12pt, a4paper]{article}

% ── Packages ──────────────────────────────────────────────────────────────────
\usepackage[utf8]{inputenc}
\usepackage[T1]{fontenc}
\usepackage[english]{babel}
\usepackage{csquotes}           % required by biblatex-apa when babel is active
\usepackage{amsmath, amssymb, amsthm}
\usepackage{mathtools}
\usepackage{geometry}
\usepackage{graphicx}
\usepackage{booktabs}
\usepackage{array}
\usepackage{multirow}
\usepackage{xcolor}
\usepackage{hyperref}
\usepackage{enumitem}
\usepackage{float}
\usepackage{caption}
\usepackage{subcaption}
\usepackage{tikz}
\usepackage{pgfplots}
\pgfplotsset{compat=1.18}
\usepackage{algorithm}
\usepackage{algpseudocode}
\usepackage[style=apa, backend=biber, natbib=true]{biblatex}
\addbibresource{references.bib}

\geometry{
  top=2.5cm, bottom=2.5cm,
  left=3cm,  right=3cm
}

\hypersetup{
  colorlinks = true,
  linkcolor  = blue!60!black,
  citecolor  = blue!60!black,
  urlcolor   = blue!60!black
}

% ── Theorem environments ──────────────────────────────────────────────────────
\theoremstyle{definition}
\newtheorem{definition}{Definition}[section]
\newtheorem{example}{Example}[section]

\theoremstyle{plain}
\newtheorem{theorem}{Theorem}[section]

% ── Custom commands ───────────────────────────────────────────────────────────
\newcommand{\R}{\mathbb{R}}
\newcommand{\Z}{\mathbb{Z}}
\newcommand{\Hk}{H_k}
\newcommand{\dgm}{\mathrm{Dgm}}
\newcommand{\pers}{\mathrm{pers}}

% ── Title ─────────────────────────────────────────────────────────────────────
\title{%
  \textbf{Topological Data Analysis for Time Series Classification}\\[0.5em]
  \large Background and Feature Design
}
\author{İsmail Güzel}
\date{May 2026}

% ──────────────────────────────────────────────────────────────────────────────
\begin{document}
\maketitle
\tableofcontents
\newpage

% ══════════════════════════════════════════════════════════════════════════════
\section{Introduction}
% ══════════════════════════════════════════════════════════════════════════════

\subsection*{Motivation: what shape does a time series have?}

Consider a time series that records, say, the hourly temperature over a month.
A statistician would summarise it with its mean, variance, autocorrelation
coefficients, and spectral density. These are all perfectly reasonable — but
they are all \emph{coordinate-dependent}: they depend on the absolute scale and
the temporal ordering of the values.

Now ask a different question: if you drew the temperature curve on paper, what
would you notice about its \emph{shape}? You might say: ``it goes up and down
periodically'', ``there is a persistent plateau around day 10'', ``there is one
very sharp spike on day 22''. These are \emph{topological} observations — they
describe how the curve connects, bends, and self-organises, independently of the
exact numerical values.

\textbf{Topological Data Analysis (TDA)} is the branch of applied mathematics
that makes these observations precise and computable
\citep{carlsson2009topology, edelsbrunner2010computational}.
The key idea is to \emph{vary a threshold} and track how connected pieces of the
data are born, merge, and disappear. The result is a concise multi-scale summary
called a \emph{persistence diagram}.

\subsection*{Why persistence diagrams are useful for time series}

The central object of TDA is the \emph{persistence diagram}, a compact summary of
how topological features are born and die as a filtration threshold is varied.
Three properties make it particularly useful in practice:

\begin{enumerate}[nosep]
  \item \textbf{Stability.} Persistence diagrams are stable under small
    perturbations: adding a little noise to a time series shifts the diagram by
    at most the same amount in the bottleneck distance
    \citep{cohen-steiner2007stability}. This is the topological analog of
    Lipschitz continuity.

  \item \textbf{Scale invariance.} The diagram captures structure that is present
    across \emph{multiple scales} simultaneously. A trend and a superimposed
    oscillation both leave traces in the same diagram.

  \item \textbf{Complementarity.} Persistence diagrams encode global shape rather
    than local statistics. They are therefore complementary to statistical feature
    extraction methods such as TSFresh \citep{christ2018time}, which excel at
    local temporal patterns, autocorrelation structure, and spectral content.
\end{enumerate}

This complementarity motivates the hybrid pipeline studied in this work: combining
18 topological features with 100 TSFresh features for the classification of
synthetic time series patterns.

\subsection*{Structure of this report}

\bigskip
\noindent
Section~\ref{sec:persistent_homology} introduces persistent homology, the
mathematical foundation, with an intuitive water-level analogy before the formal
definitions.
Section~\ref{sec:filtrations} describes the three filtration strategies applied to
univariate time series.
Section~\ref{sec:feature_extraction} details the scalar features derived from each
persistence diagram.
Section~\ref{sec:feature_selection} explains the mutual information-based selection
procedure used to reduce the feature set.
Section~\ref{sec:tda_literature} reviews relevant literature, from foundational
papers to recent (2022--2026) time series applications.
Section~\ref{sec:results_summary} summarizes the empirical performance of the
topological features on the 10-class classification benchmark used in this work.

% ══════════════════════════════════════════════════════════════════════════════
\section{Persistent Homology}
\label{sec:persistent_homology}
% ══════════════════════════════════════════════════════════════════════════════

\subsection{Intuition: the rising-water analogy}

Before introducing the formal machinery, it helps to build intuition with a
concrete analogy.

Imagine a hilly landscape and a rising sea level. As the water rises:
\begin{itemize}[nosep]
  \item First, the highest peaks emerge as isolated \emph{islands} (connected
    components are born).
  \item As the water rises further, nearby islands merge into \emph{archipelagos}
    (components die by merging with older ones).
  \item Enclosed \emph{lakes} may appear and then fill in (loops are born and die).
\end{itemize}
The \emph{birth height} of an island is the elevation of its highest peak; its
\emph{death height} is the elevation at which it merges with a taller island.
The difference — the \emph{persistence} — measures how prominent the peak is.
A tall mountain has high persistence; a small bump has low persistence.

For a time series, the ``landscape'' is the graph of the signal values. The
``water level'' is a varying threshold. As we raise (or lower) this threshold,
connected groups of time steps appear and merge, and the birth/death pairs encode
the structural features of the signal.

\subsection{Simplicial Complexes and Homology}

Let $X$ be a finite point cloud or a function defined on a topological space.
Algebraic topology studies $X$ through the lens of \emph{homology groups}
$H_k(X; \mathbb{F})$ over a field $\mathbb{F}$, where the integer $k \geq 0$
is the \emph{homological dimension}:

\begin{itemize}[nosep]
  \item $H_0$: connected components (\emph{islands} in the water analogy)
  \item $H_1$: one-dimensional loops (\emph{tunnels} through the terrain)
  \item $H_2$: two-dimensional voids (\emph{cavities} enclosed in 3D)
\end{itemize}

In practice, $X$ is replaced by a \emph{simplicial complex} built from the data,
and homology is computed combinatorially via chain complexes and boundary operators
\citep{hatcher2002algebraic}.
For a 1-dimensional time series the simplicial complex is simply a graph whose
nodes are the time steps and whose edges connect consecutive steps.

\subsection{Filtrations}

A \emph{filtration} is a nested sequence of simplicial complexes
\[
  \emptyset = K_0 \subseteq K_1 \subseteq \cdots \subseteq K_n = K
\]
parameterised by a threshold $t_0 \leq t_1 \leq \cdots \leq t_n$.
As the threshold increases, new simplices are added and the topology changes:
connected components merge, loops appear or fill in.

\begin{definition}[Birth and Death]
A topological feature (e.g., a connected component) is \emph{born} at threshold
$b$ when it first appears in the filtration and \emph{dies} at threshold $d > b$
when it merges with an older feature or is filled by a higher-dimensional simplex.
Its \emph{persistence} is $\pers = d - b$.
\end{definition}

\subsection{Persistence Diagrams}

The \emph{persistence diagram} $\dgm_k(f)$ for homological dimension $k$ is the
multiset of points $(b_i, d_i) \in \R^2$ corresponding to all features that are born
and die during the filtration, together with the diagonal $\{(t, t) : t \in \R\}$
counted with infinite multiplicity \citep{edelsbrunner2000topological}.

Features near the diagonal (small persistence $d - b$) are considered topological
noise; features far from the diagonal encode robust structural properties of the
signal.

A fundamental stability result guarantees that small perturbations of the input
function produce small changes in the persistence diagram under the bottleneck distance
$d_B$ \citep{cohen-steiner2007stability}:

\begin{theorem}[Stability of Persistence Diagrams]
Let $f, g : X \to \R$ be continuous functions on a compact topological space $X$.
Then
\[
  d_B\bigl(\dgm(f),\, \dgm(g)\bigr) \;\leq\; \|f - g\|_\infty.
\]
\end{theorem}

This stability property is crucial for applications to noisy time series: small
measurement errors cannot drastically alter the topological summary.

% ══════════════════════════════════════════════════════════════════════════════
\section{Filtrations for Univariate Time Series}
\label{sec:filtrations}
% ══════════════════════════════════════════════════════════════════════════════

Given a univariate time series $f : \{0, 1, \ldots, T-1\} \to \R$, we apply
three complementary filtration strategies. Together they form the
\texttt{sublevel+h1} pipeline used in this work.

\subsection{Sublevel Set Filtration ($\text{sub\_H0}$)}

\begin{definition}[Sublevel Set]
The \emph{sublevel set} of $f$ at threshold $t$ is
\[
  \mathcal{L}^-(t) = \{s : f(s) \leq t\}.
\]
\end{definition}

As $t$ increases from $\min f$ to $\max f$, time steps are added to the active
set in increasing order of their function value. Two active time steps $s_1, s_2$
are connected if they are adjacent in time. Connected components are born when
isolated time steps are added and die when they merge with an older component (by
the elder rule).

The resulting $H_0$ persistence diagram encodes how the \emph{local minima} of
$f$ are organised. A time series with a single dominant minimum (e.g., a
deterministic trend reaching its nadir at one point) produces a diagram with one
highly persistent point far from the diagonal, while a noisy stationary process
produces many low-persistence points clustered near the diagonal (noise). The
sublevel filtration is equivalent to computing the zeroth persistent homology of
the 1-dimensional CW complex built on the graph of $f$.

\begin{example}[Sublevel filtration on a mean-shift series]
Consider a time series $f$ with a mean of $0$ for $t < 500$ and a mean of $2$
for $t \geq 500$ (a mean shift). As the threshold $t$ rises from $-\infty$:
\begin{enumerate}[nosep, label=\roman*.]
  \item Near $t = 0$: the segment around the first half is activated; one
    connected component is born.
  \item Near $t = 2$: the segment around the second half is activated; a second
    component is born.
  \item At the threshold where the two segments touch (at the shift boundary):
    the younger component merges into the older one — it dies.
\end{enumerate}
The resulting persistence diagram contains a point at approximately $(2, 2+\delta)$
— i.e., it has low persistence — while the first component persists until
$t = \max f$, giving a high-persistence point. This asymmetric diagram is the
topological signature of a mean shift.
\end{example}

\subsection{Superlevel Set Filtration ($\text{sup\_H0}$)}

\begin{definition}[Superlevel Set]
The \emph{superlevel set} of $f$ at threshold $t$ is
\[
  \mathcal{L}^+(t) = \{s : f(s) \geq t\}.
\]
\end{definition}

As $t$ decreases from $\max f$ to $\min f$, time steps are added to the active
set in decreasing order of their function value. This filtration captures the
organisation of \emph{local maxima}: a sharp spike (point anomaly) creates a
single highly persistent component because the spike maximum appears early and
persists until the threshold drops to the baseline level. Mean shifts create step
changes in the superlevel diagram. The superlevel diagram is the mirror image of
the sublevel diagram when $f$ is replaced by $-f$.

\subsection{Takens Embedding and Vietoris--Rips Filtration ($\text{H1}$)}

The sublevel and superlevel filtrations only capture $H_0$ features (connected
components). To access $H_1$ features (loops), we apply a different approach
based on \emph{delay embedding} \citep{takens1981detecting}.

\begin{definition}[Takens Delay Embedding]
Given a time series $f$ of length $T$, a delay $\tau > 0$, and an embedding
dimension $d$, the \emph{Takens embedding} maps each time step $s$ to the point
\[
  \phi(s) = \bigl(f(s),\, f(s+\tau),\, \ldots,\, f(s+(d-1)\tau)\bigr) \in \R^d.
\]
\end{definition}

By Takens' embedding theorem \citep{takens1981detecting}, if $f$ arises from a
deterministic dynamical system, then $\phi$ reconstructs the attractor of the
system up to diffeomorphism (under generic conditions). For a periodic or
quasi-periodic signal, the attractor is a torus or circle, and $H_1$ detects the
resulting loops.

\begin{example}[Periodicity via $H_1$]
A sinusoidal signal $f(t) = \sin(2\pi t / P)$ with period $P$, when delay-embedded
with $\tau = P/4$ and $d = 2$, produces the point cloud
$\{(\sin\theta,\, \cos\theta)\}$ — a perfect circle. Under Vietoris--Rips
filtration, this circle creates exactly one persistent $H_1$ class born at small
$\epsilon$ and dying only when the circle is ``filled in'' at large $\epsilon$.
Its persistence is proportional to the signal amplitude. A stationary noise
process, by contrast, produces a point cloud without circular structure, giving
only short-lived $H_1$ classes near the diagonal.
\end{example}

The point cloud $\{\phi(s)\}$ is filtered by the \emph{Vietoris--Rips
complex}: at radius $\epsilon$, any set of points mutually within distance
$\epsilon$ forms a simplex:
\[
  \mathcal{R}_\epsilon = \bigl\{\sigma \subseteq P :
    \mathrm{diam}(\sigma) \leq \epsilon \bigr\}.
\]
As $\epsilon$ increases from 0 to $\infty$, $H_1$ classes (loops) are born and
die; a persistent loop indicates genuine cyclic or oscillatory structure in the
time series.

\medskip\noindent
\textbf{Important limitation.} The sublevel and superlevel filtrations are
constructed from the \emph{values} of the time series and are therefore invariant
to any permutation of time steps that preserves the multiset of values. They
encode ``what heights appear'' but not ``when they appear''. This means that
stationarity, autocorrelation structure, and volatility clustering — properties
that are inherently temporal — are not directly captured. The Takens embedding
partially addresses this by introducing temporal ordering through the delay
parameter $\tau$, but the basic sublevel/superlevel features share this limitation.
This is the primary reason why topological features underperform statistical
features on the \texttt{stationary} class in the experiments described in
Section~\ref{sec:results_summary}.

\bigskip
\noindent
\textbf{Why three filtrations?} Each filtration answers a different topological
question about the time series:

\begin{table}[H]
\centering
\caption{Complementary information captured by the three filtration strategies.}
\label{tab:filtrations}
\begin{tabular}{llll}
\toprule
\textbf{Diagram} & \textbf{Filtration} & \textbf{Dimension} & \textbf{Captures} \\
\midrule
$\text{sub\_H0}$ & Sublevel sets      & $H_0$ & Local minima, connectivity from below \\
$\text{sup\_H0}$ & Superlevel sets    & $H_0$ & Local maxima, connectivity from above \\
$\text{H1}$      & Takens + Rips      & $H_1$ & Cyclic structure, delay-space loops \\
\bottomrule
\end{tabular}
\end{table}

% ══════════════════════════════════════════════════════════════════════════════
\section{Scalar Feature Extraction from Persistence Diagrams}
\label{sec:feature_extraction}
% ══════════════════════════════════════════════════════════════════════════════

A persistence diagram is a \emph{multiset} of points in $\R^2$ and cannot be
directly fed into standard machine learning pipelines. Several vectorisation
methods have been proposed to convert diagrams to fixed-length feature vectors.
We use a combination of Carlsson coordinates, persistence entropy, and persistence
landscapes.

\subsection{Carlsson Coordinates}

\citet{adcock2016ring} proposed a set of polynomial functions on the space of
persistence diagrams that are invariant under rearrangement of the points.
Let $D = \{(b_i, d_i)\}$ be a persistence diagram (diagonal points excluded)
and let $p_i = d_i - b_i$ denote the persistence of the $i$-th feature.

The five Carlsson-type invariants used in this work are:

\begin{align}
  f_1(D) &= \sum_i b_i \cdot p_i                        \label{eq:carl_f1} \\[4pt]
  f_2(D) &= \sum_i (d_{\max} - d_i) \cdot p_i           \label{eq:carl_f2} \\[4pt]
  f_3(D) &= \sum_i b_i^2 \cdot p_i^4                    \label{eq:carl_f3} \\[4pt]
  f_4(D) &= \sum_i (d_{\max} - d_i)^2 \cdot p_i^4       \label{eq:carl_f4} \\[4pt]
  f_5(D) &= \max_i p_i                                  \label{eq:carl_f5_max}
\end{align}

where $d_{\max} = \max_i d_i$. The first two coordinates are linear in
persistence and use birth/death as weights — $f_1$ is the
\emph{birth-weighted total persistence} (a signed quantity that captures
\emph{where} on the filtration axis the bars are concentrated; it can be
negative when births lie below zero), and $f_2$ is the analogous
\emph{death-distance-weighted total persistence}. The third and fourth
coordinates raise both factors to higher powers ($b_i^2$ or $(d_{\max} - d_i)^2$
multiplied by $p_i^4$); the quartic in $p_i$ aggressively down-weights
low-persistence (noise) bars and emphasises long-lived (signal) features,
while the quadratic birth/death weight sharpens the contrast between bars at
extreme filtration levels and those near the centre. Finally, $f_5$ is simply
the single most persistent bar — a robust summary of the dominant topological
feature in the diagram. This particular five-statistic set follows the form
used by \citet{umeda2017time} for time-series classification rather than the
original ring-of-polynomial-invariants construction of \citet{adcock2016ring},
which proposes a larger family of polynomial features.

\subsection{Persistent Entropy}

Persistent entropy \citep{chintakunta2015entropy} adapts the Shannon entropy
formula to measure the spread of persistence values in a diagram:

\begin{equation}
  E(D) = -\sum_i \hat{p}_i \log \hat{p}_i,
  \qquad \hat{p}_i = \frac{p_i}{\sum_j p_j},
  \label{eq:pers_entropy}
\end{equation}

where the sum runs over all off-diagonal points. $E(D) = 0$ when the diagram
has a single point of maximal persistence (completely concentrated signal),
and $E(D)$ is maximal when all persistence values are equal (uniform diagram,
typical of white noise). Persistent entropy thus functions as a complexity
measure: low entropy indicates a dominant structural feature; high entropy
indicates homogeneous noise.

\subsection{Persistence Landscapes}

The persistence landscape \citep{bubenik2015statistical} is a functional
representation of the persistence diagram that maps each point $(b_i, d_i)$
to a tent function:
\[
  \lambda_i(t) = \begin{cases}
    t - b_i & \text{if } b_i \leq t \leq \tfrac{b_i+d_i}{2}, \\
    d_i - t & \text{if } \tfrac{b_i+d_i}{2} \leq t \leq d_i, \\
    0 & \text{otherwise.}
  \end{cases}
\]
The $k$-th landscape function $\lambda_k(t)$ is the $k$-th largest value among
all $\lambda_i(t)$ at each threshold $t$. The landscape lives in a Banach space
and admits a well-defined mean and variance over samples, making it amenable to
statistical inference.

\medskip\noindent
\textbf{Discretisation used in this work.} We discretise the filtration axis
into $R = 100$ uniform grid points spanning $[\min_i b_i,\, \max_i d_i]$, and
on each grid point we store the top-$K$ landscape values
$\lambda_1(t),\,\ldots,\,\lambda_K(t)$ with $K = 5$. The five landscape levels
are then averaged into a single discrete function
\begin{equation}
  \bar\lambda(t) \;=\; \frac{1}{K} \sum_{k=1}^{K} \lambda_k(t),
\end{equation}
which we treat as a vector $\bar\lambda \in \R^R$. Two scalar summaries are
extracted from $\bar\lambda$:
\begin{equation}
  \ell_1 \;=\; \|\bar\lambda\|_1 \;=\; \sum_{r=1}^{R} |\bar\lambda(t_r)|,
  \qquad
  \ell_2 \;=\; \|\bar\lambda\|_2 \;=\; \sqrt{\sum_{r=1}^{R} \bar\lambda(t_r)^2}.
  \label{eq:landscapes}
\end{equation}

The $L^1$ norm captures the total ``area'' of the averaged landscape — large
when many bars overlap in the same region of the filtration axis. The $L^2$
norm is closer to a peak detector — large when one or two regions dominate.
Together they summarise the dominance hierarchy of topological features while
remaining computationally cheap (no per-level vectorisation is stored).

\subsection{Summary: 8 Features per Diagram}

\begin{table}[H]
\centering
\caption{Scalar features extracted from each persistence diagram.}
\label{tab:features_per_diagram}
\begin{tabular}{clll}
\toprule
\textbf{Index} & \textbf{Feature} & \textbf{Equation} & \textbf{Interpretation} \\
\midrule
1 & \texttt{carl\_f1}       & \eqref{eq:carl_f1}      & Birth-weighted total persistence (signed) \\
2 & \texttt{carl\_f2}       & \eqref{eq:carl_f2}      & Death-distance-weighted total persistence \\
3 & \texttt{carl\_f3}       & \eqref{eq:carl_f3}      & $b^2 \times p^4$: long bars at extreme births \\
4 & \texttt{carl\_f4}       & \eqref{eq:carl_f4}      & $(d_{\max}-d)^2 \times p^4$: long bars dying early \\
5 & \texttt{carl\_f5\_max}  & \eqref{eq:carl_f5_max}  & Most persistent bar (single max) \\
6 & \texttt{entropy}        & \eqref{eq:pers_entropy} & Complexity / spread of persistence values \\
7 & \texttt{landscape\_l1}  & \eqref{eq:landscapes}   & $L^1$ norm of mean persistence landscape \\
8 & \texttt{landscape\_l2}  & \eqref{eq:landscapes}   & $L^2$ norm of mean persistence landscape \\
\bottomrule
\end{tabular}
\end{table}

\noindent
Applying these 8 features to each of the 3 diagrams yields $3 \times 8 = 24$ raw
topological features per time series. After mutual information-based feature selection
(Section~\ref{sec:feature_selection}), 18 of these are retained.

The full feature naming convention is:
\[
  \texttt{topo\_\_\{diagram\}\_\_\{feature\}}
\]
where \texttt{diagram} $\in \{\texttt{sub\_H0},\, \texttt{sup\_H0},\, \texttt{H1}\}$
and \texttt{feature} $\in \{\texttt{carl\_f1},\, \texttt{carl\_f2},\, \texttt{carl\_f3},\,
\texttt{carl\_f4},\, \texttt{carl\_f5\_max},\, \texttt{entropy},\,
\texttt{landscape\_l1},\, \texttt{landscape\_l2}\}$.

% ══════════════════════════════════════════════════════════════════════════════
\section{Feature Selection}
\label{sec:feature_selection}
% ══════════════════════════════════════════════════════════════════════════════

Feature selection is performed using \emph{mutual information} (MI), a
non-parametric measure of statistical dependence:
\begin{equation}
  I(X; Y) = \sum_{x,y} p(x, y) \log \frac{p(x,y)}{p(x)\,p(y)}.
  \label{eq:mi}
\end{equation}
MI equals zero if and only if $X$ and $Y$ are statistically independent, and
is invariant under monotone transformations. It captures both linear and
non-linear dependencies, making it more general than Pearson correlation for
the highly non-Gaussian distribution of persistence-based features.

\medskip
\noindent
\textbf{Leakage prevention.} MI is estimated on the \emph{training split only}
and then applied to select features before any model sees the test data. Computing
MI on the full dataset would leak test-label information into feature selection and
inflate reported accuracy \citep{kaufman2012leakage}. In practice, MI is estimated
using the $k$-nearest-neighbour estimator of \citet{kraskov2004estimating} as
implemented in \texttt{sklearn.feature\_selection.mutual\_info\_classif}
\citep{pedregosa2011scikit}.

\medskip
\noindent
\textbf{Selection result.} Of the 24 raw topological features, 18 are retained
on the 10-class training split. The 6 dropped features are
\texttt{carl\_f5\_max} and \texttt{landscape\_l2} from each of the three
diagrams (\texttt{H1}, \texttt{sub\_H0}, \texttt{sup\_H0}).

\medskip\noindent
\textbf{Why high-MI features are dropped.} It would be incorrect to attribute
these removals solely to low MI: in fact, \texttt{sup\_H0\_\_carl\_f5\_max}
and \texttt{sup\_H0\_\_landscape\_l2} carry some of the \emph{highest} MI scores
in the pool. The selection procedure combines MI ranking with a
\emph{correlation-based redundancy filter}: among groups of features whose pairwise
Pearson correlation exceeds a threshold, only one representative is kept. The
six dropped features fall into two structurally redundant pairs per diagram:
\begin{itemize}[nosep]
  \item \texttt{carl\_f5\_max} (maximum bar lifetime) is highly correlated with
    \texttt{landscape\_l1} (area under the mean landscape), because both grow
    monotonically with the length of the dominant persistent bar;
  \item \texttt{landscape\_l2} ($L^2$ norm of the mean landscape) is highly
    correlated with \texttt{landscape\_l1} ($L^1$ norm of the same vector) since
    both are norms of the same discrete function.
\end{itemize}
The retained 18 features therefore represent the same topological information
in a less redundant form, rather than a strictly higher-MI subset.

% ══════════════════════════════════════════════════════════════════════════════
\section{Literature Review: TDA for Time Series}
\label{sec:tda_literature}
% ══════════════════════════════════════════════════════════════════════════════

\subsection{Foundations}

Persistent homology was formalised in a computational context by
\citet{edelsbrunner2000topological} and
\citet{zomorodian2005computing}, who introduced efficient algorithms for its
computation. The stability theorem --- critical for applications to noisy data ---
was proved by \citet{cohen-steiner2007stability}.
\citet{carlsson2009topology} provided the first broad survey of TDA as a practical
data analysis methodology, establishing the language and goals that shaped subsequent
applied work.

\subsection{Vectorisation Methods}

Converting persistence diagrams to machine-learning-compatible vectors has been an
active research area. \citet{bubenik2015statistical} introduced persistence
landscapes, establishing their statistical properties (mean, variance, central limit
theorem) in a Hilbert space setting. \citet{adcock2016ring} characterised the ring
of polynomial invariants of persistence diagrams, motivating the Carlsson coordinates
used in this work. \citet{chintakunta2015entropy} proposed persistent entropy as a
single-scalar complexity measure. More recently, persistence images
\citep{adams2017persistence} and the sliced Wasserstein kernel
\citep{carriere2017sliced} have extended the vectorisation toolkit; however, for
the present application the scalar features described above are computationally
efficient and interpretable.

\subsection{TDA Applied to Time Series}

\citet{perea2015sliding} introduced the sliding window embedding (a variant of
Takens embedding) and proved that it produces circles in the delay space for
periodic signals, providing a rigorous foundation for using $H_1$ features to
detect periodicity. The period manifests as a highly persistent $H_1$ class with
$\pers \approx$ (signal amplitude).

\citet{umeda2017time} applied sublevel and superlevel filtrations directly to
financial and physiological time series, showing that persistence diagrams
discriminate between different regimes more robustly than classical spectral
features. \citet{gidea2018topological} demonstrated that the $L^1$ norm of the
first persistence landscape detects early warning signals before financial crashes
— a structural break detection task directly analogous to the \texttt{trend\_shift}
and \texttt{variance\_shift} classes studied here.

\citet{seversky2016time} performed a systematic comparison of TDA-based features
against classical time series descriptors on UCR classification benchmarks and
found that sublevel persistence features compete with or outperform dynamic time
warping and shapelet methods on datasets with irregular structural transitions,
while underperforming on noise-heavy or purely stochastic datasets.

More recent work has pushed TDA for time series in directions that are directly
relevant to the present pipeline. \citet{lin2022temporal} proposed a
\emph{temporal filtration} that augments standard persistence with explicit time
ordering information, addressing a key limitation of purely value-based sublevel
and superlevel constructions. \citet{elyaagoubi2023multivariate} extended the
discussion from univariate trajectories to multivariate time series, especially
brain-signal settings where both channel-wise evolution and cross-channel
structure matter. \citet{ferra2023importance} combined persistence landscapes of
time series with neural networks and an attribution layer, showing that
topological summaries can be both predictive and interpretable in supervised
classification settings.

Bridging the gap between theoretical tools and industrial anomaly detection,
\citet{perez2021topological} applied TDA to multivariate sensor time series from
manufacturing processes, showing that persistence-based features detect hidden
structural patterns that remain invisible to classical multivariate statistics.
This is especially relevant to the anomaly classes (\texttt{collective\_anomaly},
\texttt{point\_anomaly}, \texttt{variance\_shift}) studied in the present work.

Application-oriented studies from 2024 further support the usefulness of TDA for
complex sequential data. \citet{flammer2024chaotic} used persistent homology for
classification of chaotic multivariate EEG windows and reported that topology is
effective for detecting seizure-related structure after suitable dimension
reduction. \citet{zheng2024multivariate} proposed a multichannel EEG pipeline
that couples Hilbert--Huang signal representations with topological features,
illustrating how TDA can be integrated with more classical signal-processing
front ends rather than used in isolation.

\citet{chazal2021introduction} provides a comprehensive modern introduction to
TDA with statistical perspectives.

\subsection{TDA and Machine Learning}

The combination of TDA features with standard machine learning classifiers has
been explored widely in recent years. \citet{carlsson2020topological} provides
an authoritative review of topological pattern recognition for point cloud data,
describing the theoretical framework within which persistence-based classifiers
operate. At the applied level, \citet{hensel2021survey} surveyed machine learning
applications of TDA across domains, noting that the combination of
persistence-based features with gradient-boosted trees is the most consistently
competitive approach, with random forests close behind — consistent with the
CatBoost and Random Forest results reported in Section~\ref{sec:results_summary}.

\citet{giotto2021} describe the \texttt{giotto-tda} Python library, which implements
the persistent homology pipeline used in this work, including Vietoris--Rips
computation for Takens embeddings and sublevel/superlevel filtrations for scalar
functions. The underlying computation of Vietoris--Rips persistent homology is
performed by the \textsc{Ripser} algorithm \citep{bauer2021ripser}, one of the
fastest available implementations; a Python interface is provided by
\citet{tralie2018ripser}.

For broader context, \citet{su2025review} review recent developments in TDA and
topological deep learning beyond classical persistent homology. Although their
scope is wider than time series alone, the review is useful for situating modern
sequence-analysis pipelines within the rapidly expanding ecosystem of topological
representations, vectorisations, and learning architectures.

% ══════════════════════════════════════════════════════════════════════════════
\section{Empirical Performance Summary}
\label{sec:results_summary}
% ══════════════════════════════════════════════════════════════════════════════

The topological features were evaluated on a 10-class synthetic time series
classification task (\texttt{topo-test} dataset: 1,000 series, 100 per class,
length 1,000 time points). Results are compared against a 100-feature TSFresh
pipeline and a 118-feature hybrid concatenation. All experiments use an 80/20
stratified train/test split with seed 42. Table~\ref{tab:results_main} summarises
test accuracy for four classifiers across all three feature pipelines.

\begin{table}[H]
\centering
\caption{Test accuracy on the 10-class topo-test benchmark.}
\label{tab:results_main}
\begin{tabular}{lccccc}
\toprule
\textbf{Pipeline} & \textbf{Features} & \textbf{RF} & \textbf{XGBoost}
  & \textbf{CatBoost} & \textbf{SVM} \\
\midrule
TSFresh           & 100 & 93.0\% & 94.5\% & \textbf{95.0\%} & 89.0\% \\
\textbf{Topology} & \textbf{18}  & 85.5\% & 85.0\% & 86.5\% & 78.5\% \\
Hybrid            & 118 & 92.5\% & 94.5\% & 94.5\% & \textbf{91.0\%} \\
\bottomrule
\end{tabular}
\end{table}

\noindent
\textbf{Feature efficiency.} The topological pipeline delivers 86.5\% CatBoost
accuracy from just 18 features — using $100/18 \approx 5.6\times$ fewer features
than the TSFresh pipeline at the cost of an 8.5 percentage point accuracy
reduction. Expressed as accuracy per feature:

\begin{align*}
  \text{TSFresh:} &\quad 95.0\,\% / 100 = 0.95\,\% \text{ per feature,} \\
  \text{Topology:} &\quad 86.5\,\% / 18  = 4.81\,\% \text{ per feature.}
\end{align*}

The accuracy-per-feature ratio is therefore
$4.81 / 0.95 \approx 5.1\times$ — topological features achieve
approximately \textbf{5$\times$ better accuracy-per-feature efficiency} than
the statistical baseline.

\medskip
\noindent
\textbf{Per-class highlights.}
Table~\ref{tab:per_class} reports CatBoost F1 scores by class across all three
pipelines.

\begin{table}[H]
\centering
\caption{CatBoost test F1 scores by class across the three feature pipelines.
  Classes where topology matches or exceeds TSFresh are marked ($\dagger$).}
\label{tab:per_class}
\begin{tabular}{lccc}
\toprule
\textbf{Class} & \textbf{TSFresh} & \textbf{Topology} & \textbf{Hybrid} \\
\midrule
stationary          & 0.884 & 0.651 & 0.889 \\
deterministic\_trend & 0.974 & 0.872 & 0.974 \\
stochastic\_trend    & 1.000 & 0.895 & 0.974 \\
volatility          & 0.865 & 0.750 & 0.889 \\
collective\_anomaly  & 0.976 & 0.732 & 0.947 \\
contextual\_anomaly  & 1.000 & 0.974 & 1.000 \\
mean\_shift          & 0.909 & 0.850 & 0.870 \\
point\_anomaly       & 1.000 & 1.000$\dagger$ & 1.000$\dagger$ \\
trend\_shift         & 0.919 & \textbf{0.974}$\dagger$ & 0.919 \\
variance\_shift      & 0.974 & \textbf{0.976}$\dagger$ & 1.000$\dagger$ \\
\midrule
\textbf{Macro F1}   & \textbf{0.950} & 0.867 & 0.946 \\
\bottomrule
\end{tabular}
\end{table}

\noindent
Topology \emph{surpasses} TSFresh on \texttt{trend\_shift} (0.974 vs.\ 0.919)
and matches it on \texttt{variance\_shift} and \texttt{point\_anomaly}. Both of
these advantages are structurally interpretable:

\begin{itemize}
  \item \textbf{Trend shift:} A slope change creates a distinctive asymmetry in
    the \emph{sublevel} filtration: the segment on one side of the change point
    has a different minimum than the other, and the two sublevel components
    merge at a filtration level determined by the change-point amplitude. The
    sublevel persistence diagram therefore contains a single high-persistence
    point whose birth captures the lower of the two minima, which is exactly the
    pattern that \texttt{sub\_H0\_\_carl\_f3} emphasises (large $b^2 \cdot p^4$
    at extreme births). This is consistent with the empirical finding that
    \texttt{sub\_H0\_\_carl\_f3} is the single most important feature in the
    hybrid model. TSFresh chunk linear-trend features miss this signature
    whenever the change point falls between chunks.

  \item \textbf{Variance shift:} A step increase in variance produces a visible
    change in the density of sublevel crossings. Carlsson coordinates and persistent
    entropy in the sublevel diagram encode this directly. TSFresh change-quantile
    features capture similar information but are less robust to the exact location
    of the shift.

  \item \textbf{Stationary (worst case):} TSFresh has a decisive advantage because
    autocorrelation features, spectral flatness, and stationarity-related statistics
    directly encode the absence of a trend or structural change. Persistence diagrams,
    being invariant to temporal ordering, cannot detect stationarity without
    supplementary temporal information.
\end{itemize}

\medskip
\noindent
\textbf{Most important topological feature.}
In the hybrid CatBoost model (118 features: 100 TSFresh + 18 topological),
\texttt{topo\_\_sub\_H0\_\_carl\_f3} is the \emph{single most important feature}
with an importance score of 13.2\%, outranking all 100 TSFresh features.
This coordinate (Equation~\ref{eq:carl_f3}) takes the form
\[
  f_3(D) = \sum_i b_i^2 \cdot (d_i - b_i)^4,
\]
combining a quadratic birth weight with a quartic persistence weight. The
$b_i^2$ factor amplifies bars born at extreme filtration values, while $p_i^4$
aggressively suppresses short (noise) bars. The result is a single scalar that
fires strongly when a series contains \emph{long-lived bars whose births occur
at extreme signal levels} — exactly the pattern produced by trend-like
asymmetric global minima and persistent secondary minima of anomaly classes.
Stationary processes, by contrast, produce many short-lived
low-persistence components with births clustered near zero, so the $b^2 \cdot
p^4$ weighting drives $f_3$ toward small values. This is why
\texttt{sub\_H0\_\_carl\_f3} is highly discriminative across the ten classes
studied here.

\subsection*{Practical limitations of the topological pipeline}

For balanced use of topological features it is worth being explicit about their
limitations, several of which have been noted in the literature:

\begin{enumerate}
  \item \textbf{Temporal order invariance.} As discussed in
    Section~\ref{sec:filtrations}, sublevel and superlevel diagrams are invariant
    to permutations of the time axis. This means that autocorrelation, stationarity,
    and volatility clustering — all of which depend on temporal ordering — are not
    directly representable. The Takens embedding ($H_1$) partially restores temporal
    structure, but only for the delay parameter $\tau$ chosen.

  \item \textbf{Computational cost.} Computing Vietoris--Rips persistent homology
    scales poorly with the number of points. For a time series of length $T$ with
    embedding dimension $d$, the Takens point cloud has $T - (d-1)\tau$ points, and
    Rips computation is $O(n^3)$ in the worst case. In practice, the
    \textsc{Ripser} algorithm \citep{bauer2021ripser} dramatically reduces this cost
    via matrix reduction, and sublevel/superlevel filtrations remain $O(T \log T)$.

  \item \textbf{Hyperparameter sensitivity.} The Takens embedding requires choosing
    $\tau$ (delay) and $d$ (embedding dimension). A poorly chosen $\tau$ may produce
    a degenerate point cloud. In practice, the false nearest neighbours method or
    mutual information criterion is used to select $\tau$ \citep{takens1981detecting,
    perea2015sliding}.

  \item \textbf{Univariate focus.} The filtrations described here are designed for
    univariate series. Extending them to multivariate time series requires additional
    choices (e.g., channel-wise vs.\ joint embeddings), an active research area
    \citep{elyaagoubi2023multivariate, zheng2024multivariate}.
\end{enumerate}

% ══════════════════════════════════════════════════════════════════════════════
\section{Conclusion}
% ══════════════════════════════════════════════════════════════════════════════

Topological Data Analysis provides a principled, mathematically grounded approach
to time series feature extraction that is complementary to classical statistical
methods. The persistent homology framework, applied through sublevel, superlevel,
and Takens delay-embedding filtrations, extracts 18 scalar features that:

\begin{enumerate}
  \item Achieve 86.5\% classification accuracy on a 10-class synthetic benchmark
    with only 18 features ($\sim 5.1\times$ better accuracy-per-feature than the
    100-feature TSFresh pipeline);
  \item Outperform TSFresh on structural break detection tasks (\texttt{trend\_shift},
    \texttt{variance\_shift});
  \item Contribute the single most important feature in a 118-feature hybrid model
    (\texttt{sub\_H0\_\_carl\_f3}, importance 13.2\%);
  \item Provide stable, theoretically well-founded representations (Stability Theorem)
    that are robust to small measurement perturbations.
\end{enumerate}

The primary limitation of topological features is their insensitivity to temporal
ordering: persistence diagrams are invariant to permutation of time steps, which
makes them poor descriptors for stationarity, autocorrelation structure, and
volatility clustering — properties that TSFresh captures well. The hybrid pipeline
exploits both strengths and is recommended when both structural change detection
and temporal regularity characterisation are required.

% ══════════════════════════════════════════════════════════════════════════════
% References
% ══════════════════════════════════════════════════════════════════════════════
\newpage
\printbibliography

\end{document}
